\textbf{Phần 1. Trắc nghiệm lựa chọn}
\setcounterpageref{ex}{0}

\begin{ex} %[2D10V54]
Số liền trước của số chẵn bé nhất có ba chữ số là
\choice
{$101$}
{$100$}
{\True $99$}
{$98$}
\loigiai{
Số chẵn bé nhất có 3 chữ số là 100. Số liền trước là 99.
\textbf{Đáp án của câu hỏi là C}
}
\end{ex}

\begin{ex} %[2D10V54]
Số gồm $6$ trăm và $6$ đơn vị được viết là
\choice
{\True $606$}
{$660$}
{$66$}
{$6006$}
\loigiai{
Hàng chục bằng 0 nên viết là $606$.
\textbf{Đáp án của câu hỏi là A}
}
\end{ex}

\begin{ex} %[2D10V54]
Tổng của số lớn nhất có hai chữ số và số bé nhất có ba chữ số là
\choice
{$198$}
{$200$}
{$109$}
{\True $199$}
\loigiai{
$99 + 100 = 199$.
\textbf{Đáp án của câu hỏi là D}
}
\end{ex}

\begin{ex} %[2D10V54]
Trong các số sau, số nào gần với số $350$ nhất?
\choice
{$341$}
{\True $349$}
{$359$}
{$339$}
\loigiai{
349 cách 350 đúng 1 đơn vị, gần nhất.
\textbf{Đáp án của câu hỏi là B}
}
\end{ex}

\begin{ex} %[2D10V54]
Cho số có ba chữ số, biết chữ số hàng trăm gấp $3$ lần hàng đơn vị và hàng chục bằng $0$. Số đó có thể là
\choice
{$309$}
{$603$}
{\True $903$}
{$301$}
\loigiai{
Chữ số hàng trăm là 9, hàng đơn vị là 3, vì $9 : 3 = 3$. Số đó là 903.
\textbf{Đáp án của câu hỏi là C}
}
\end{ex}

\begin{ex} %[2D10V54]
Tìm một số biết rằng nếu lấy số đó cộng với $200$ rồi trừ đi $50$ thì được $650$.
\choice
{\True $500$}
{$400$}
{$600$}
{$450$}
\loigiai{
Số đó là: $650 + 50 - 200 = 500$.
\textbf{Đáp án của câu hỏi là A}
}
\end{ex}

\begin{ex} %[2D10V54]
Đếm thêm $5$ đơn vị từ số $245$ đến số $265$ thì được bao nhiêu số?
\choice
{$4$ số}
{$6$ số}
{$7$ số}
{\True $5$ số}
\loigiai{
Các số là: 245, 250, 255, 260, 265. Có 5 số.
\textbf{Đáp án của câu hỏi là D}
}
\end{ex}

\begin{ex} %[2D10V54]
Một rạp chiếu phim có hai phòng chiếu. Phòng một có $250$ ghế, phòng hai có $280$ ghế. Cả hai phòng có bao nhiêu ghế?
\choice
{$520$ ghế}
{\True $530$ ghế}
{$500$ ghế}
{$430$ ghế}
\loigiai{
$250 + 280 = 530$ (ghế).
\textbf{Đáp án của câu hỏi là B}
}
\end{ex}

\begin{ex} %[2D10V54]
Hình tam giác có độ dài ba cạnh đều bằng $100\text{ cm}$. Chu vi của hình tam giác đó là
\choice
{$200\text{ cm}$}
{$400\text{ cm}$}
{\True $300\text{ cm}$}
{$100\text{ cm}$}
\loigiai{
$100 \times 3 = 300$ (cm).
\textbf{Đáp án của câu hỏi là C}
}
\end{ex}

\begin{ex} %[2D10V54]
Tìm số tự nhiên có ba chữ số biết chữ số hàng trăm là $4$, nếu xóa chữ số hàng trăm thì được số $25$.
\choice
{\True $425$}
{$452$}
{$420$}
{$254$}
\loigiai{
Số đó là 425.
\textbf{Đáp án của câu hỏi là A}
}
\end{ex}

\begin{ex} %[2D10V54]
Một cửa hàng bán được $100\text{ kg}$ đường vào ngày thứ nhất. Ngày thứ hai bán gấp đôi ngày thứ nhất. Cả hai ngày bán được bao nhiêu ki-lô-gam đường?
\choice
{$200\text{ kg}$}
{$250\text{ kg}$}
{$350\text{ kg}$}
{\True $300\text{ kg}$}
\loigiai{
Ngày thứ hai: $100 \times 2 = 200$ kg. Cả hai ngày: $100 + 200 = 300$ kg.
\textbf{Đáp án của câu hỏi là D}
}
\end{ex}

\begin{ex} %[2D10V54]
Khối 2 trồng được $400$ cây, khối 3 trồng nhiều hơn khối 2 là $200$ cây. Cả hai khối trồng được tất cả bao nhiêu cây?
\choice
{$600$ cây}
{\True $1 000$ cây}
{$800$ cây}
{$900$ cây}
\loigiai{
Khối 3: $400 + 200 = 600$ cây. Cả hai khối: $400 + 600 = 1 000$ cây.
\textbf{Đáp án của câu hỏi là B}
}
\end{ex}

\textbf{Phần 2. Trắc nghiệm đúng sai}
\setcounterpageref{ex}{0}

\begin{ex} %[2D10V54]
Xét tính đúng sai của các khẳng định về số học phạm vi 1 000:
\choiceTF
{\True 1 000 là số tự nhiên có bốn chữ số nhỏ nhất.}
{Số 780 đọc là bảy trăm linh tám.}
{\True Hiệu giữa 1 000 và 1 bằng 999.}
{Số gồm 3 trăm và 4 chục là 304.}
\loigiai{
\begin{itemize}
    \item Ý a) Đúng: Đúng.
    \item Ý b) Sai: 780 đọc là bảy trăm tám mươi.
    \item Ý c) Đúng: 1 000 - 1 = 999.
    \item Ý d) Sai: 3 trăm và 4 chục là 340, không phải 304.
\end{itemize}
\begin{itemize}
    \item a) \textbf{Đúng}.
    \item b) \textbf{Sai}.
    \item c) \textbf{Đúng}.
    \item d) \textbf{Sai}.
\end{itemize}
}
\end{ex}

\begin{ex} %[2D10V54]
Cho dãy số: 100, 200, 300, 400, 500, 600, 700, 800, 900, 1 000:
\choiceTF
{\True Dãy số gồm 10 số đều là các số tròn trăm.}
{\True Mỗi số trong dãy cách số liền sau 100 đơn vị.}
{Số ở chính giữa dãy số là 500.}
{\True Tổng của số đầu và số cuối là 1 100.}
\loigiai{
\begin{itemize}
    \item Ý a) Đúng: Đúng, 10 số từ 100 đến 1 000.
    \item Ý b) Đúng: Đúng.
    \item Ý c) Sai: Dãy có 10 số nên không có một số chính giữa (ở giữa là 500 và 600).
    \item Ý d) Đúng: 100 + 1 000 = 1 100. Đúng.
\end{itemize}
\begin{itemize}
    \item a) \textbf{Đúng}.
    \item b) \textbf{Đúng}.
    \item c) \textbf{Sai}.
    \item d) \textbf{Đúng}.
\end{itemize}
}
\end{ex}

\textbf{Phần 3. Trả lời ngắn}
\setcounterpageref{ex}{0}

\begin{ex} %[2D10V54]
Tìm số có ba chữ số biết chữ số hàng trăm là 5, hàng chục gấp đôi hàng đơn vị và tổng ba chữ số bằng 11.
\shortans{542}
\loigiai{
\textbf{Đáp số:} $542$\\
Tổng hàng chục và đơn vị là $11 - 5 = 6$. Đơn vị là $6 : 3 = 2$, chục là 4: số 542.
}
\end{ex}

\begin{ex} %[2D10V54]
Bác An có 800 quả trứng, bán đi 500 quả. Số còn lại chia đều vào 3 thùng. Mỗi thùng có bao nhiêu quả trứng?
\shortans{100}
\loigiai{
\textbf{Đáp số:} $100$\\
Còn lại: $800 - 500 = 300$ quả. Mỗi thùng: $300 : 3 = 100$ quả.
}
\end{ex}

\begin{ex} %[2D10V54]
Có bao nhiêu số có ba chữ số mà hàng trăm là 7 và hàng đơn vị là 5?
\shortans{10}
\loigiai{
\textbf{Đáp số:} $10$\\
Dạng $7a5$ với $a \in \{0, 1, \dots, 9\}$. Có 10 số.
}
\end{ex}

\begin{ex} %[2D10V54]
Tính tổng của số chẵn lớn nhất có ba chữ số và số lẻ bé nhất có một chữ số.
\shortans{999}
\loigiai{
\textbf{Đáp số:} $999$\\
998 + 1 = 999.
}
\end{ex}

\begin{ex} %[2D10V54]
Một đội xe có 3 xe chở hàng, mỗi xe chở 300 bao xi măng. Đội xe chở tất cả bao nhiêu bao xi măng?
\shortans{900}
\loigiai{
\textbf{Đáp số:} $900$\\
3 x 300 = 900 bao.
}
\end{ex}

\begin{ex} %[2D10V54]
Cho số 654. Nếu đổi chỗ chữ số hàng trăm và chữ số hàng đơn vị thì được số mới là bao nhiêu?
\shortans{456}
\loigiai{
\textbf{Đáp số:} $456$\\
Đổi 6 và 4 ta được 456.
}
\end{ex}
